% !TeX program = xelatex
\documentclass[aspectratio=169]{beamer}

% The theme file beamerthemehitspbu.sty must be in the same folder as this tex.
\usetheme{hitspbu}
\usefonttheme{professionalfonts}

% Common packages for mathematical and technical presentations.
\usepackage{amsmath, amssymb, amsfonts}
\usepackage{mathtools}
\usepackage{bm}
\usepackage{booktabs, tabularx, array}
\usepackage{tikz}
\usetikzlibrary{positioning, arrows.meta}
\usepackage{enumitem}
\setlist[itemize]{left=0pt..1em}

% Optional helper macros.
\newcommand{\RR}{\mathbb{R}}
\newcommand{\E}{\mathbb{E}}
\newcolumntype{Y}{>{\raggedright\arraybackslash}X}

% Metadata. Replace these placeholders with your own information.
\title[Short Presentation Title]{Full Presentation Title\\Second Title Line If Needed}
\author[Presenter]{Presenter Name}
\renewcommand{\hitsupervisor}{Supervisor Name}
\institute[Short Institute]{Department or School, University Name}
\date{Month Day, Year}

% Optional title-page controls. Uncomment and adjust when needed.
% \setlength{\hitTitlePageTitleY}{1.40cm}
% \setlength{\hitTitlePageAuthorX}{-0.92cm}
% \setlength{\hitTitlePageAuthorY}{-0.82cm}
% \setlength{\hitTitlePageSupervisorY}{-1.42cm}
% \renewcommand{\hitTitlePageAuthorAlign}{right}
% \renewcommand{\hitTitlePageAuthorAnchor}{east}
% \renewcommand{\hitTitlePageAuthorBase}{current page.east}

% Optional body-page title controls.
% \setlength{\hitFrameTitleLeft}{5.10cm}
% \setlength{\hitFrameTitleRight}{2.20cm}
% \setlength{\hitFrameTitleY}{-0.66cm}
% \renewcommand{\hitFrameTitleFont}{\hitFrameTitleFace\fontsize{20}{24}\selectfont}

\begin{document}

\begin{frame}
  \titlepage
\end{frame}

\begin{frame}{Contents}
  \tableofcontents
\end{frame}

\section{Basic Content}

\begin{frame}{Itemize and Blocks}
\begin{block}{Standard Block}
Use the usual Beamer environments. The first-level item marker is a small hollow circle.
\end{block}

\begin{itemize}
  \item First point with a concise statement.
  \item Second point with a short explanation.
  \item Third point with a displayed equation:
  \[
    \min_{x\in\RR^d} f(x)+g(Bx).
  \]
\end{itemize}

\begin{alertblock}{Alert Block}
Use alert blocks for assumptions, warnings, or key conditions.
\end{alertblock}
\end{frame}

\begin{frame}{Columns}
\begin{columns}[T,onlytextwidth]
\column{0.48\textwidth}
\begin{block}{Left Column}
\begin{itemize}
  \item Model setup.
  \item Main assumptions.
  \item Notation.
\end{itemize}
\end{block}

\column{0.48\textwidth}
\begin{exampleblock}{Right Column}
\[
  \E\|x_k-x^\star\|^2 \le C\rho^k.
\]
This block is useful for conclusions or examples.
\end{exampleblock}
\end{columns}
\end{frame}

\section{Structured Slides}

\begin{frame}{Table}
\centering
\begin{tabularx}{0.78\textwidth}{@{}lYY@{}}
\toprule
Item & Description & Remark \\
\midrule
Method A & Deterministic baseline & Stable \\
Method B & Stochastic variant & Scalable \\
Method C & Accelerated method & Fast \\
\bottomrule
\end{tabularx}
\end{frame}

\begin{frame}{Figure Placeholder}
\centering
\begin{tikzpicture}
  \draw[HITRed, line width=1pt, rounded corners=3pt] (0,0) rectangle (8.6,3.8);
  \node[align=center, text=HITSlate] at (4.3,1.9)
    {Replace this TikZ placeholder with\\\texttt{\textbackslash includegraphics[width=...]\{file-name\}}};
\end{tikzpicture}
\end{frame}

\begin{frame}{Theorem Style}
\begin{theorem}
Assume that $f$ is convex and $L$-smooth. Then a suitable first-order method satisfies
\[
  F(x_k)-F(x^\star) \le \mathcal{O}\left(\frac{1}{k}\right).
\]
\end{theorem}
\begin{proof}
State only the main proof idea on slides; keep technical details in backup pages.
\end{proof}
\end{frame}

\section{Special Frames}

\begin{frame}[plain]{Plain Frame}
\centering
This frame intentionally has no header and no footer.

\vspace{1em}
Use it for full-page figures, section breaks, or special transition pages.
\end{frame}

\begin{hitfullframe}
\centering
The custom \texttt{hitfullframe} environment is a shortcut for a plain full-content frame.
\end{hitfullframe}

\begin{frame}{Normal Frame After Plain}
\begin{itemize}
  \item The normal header, footer, and page number return automatically.
  \item Use standard \texttt{frame} environments for most slides.
\end{itemize}
\end{frame}

\appendix

\begin{frame}[noframenumbering]{Backup Slide Not Counted}
This slide keeps the template style, but it is not included in the total frame count.

\[
  \text{Displayed page numbers remain like } 1/20,\;2/20,\ldots,20/20.
\]
\end{frame}

\begin{frame}[plain,noframenumbering]{Plain Backup Not Counted}
\centering
This backup slide has no header/footer and is also not counted.
\end{frame}

\end{document}
